\chapter{Notation and Source-Manuscript Bridges}
\label{app:notation-bridges}
\small
\sloppy
\hbadness=10000

\section{Inspected source identity}

The retained-space, Bloch-frame, and aligned-difference bridges in this appendix were
prepared from the clean \texttt{ksdft2effmass} worktree on branch
\decl{work/monograph-representation-carriers-and-derivation-graphs} at commit
\decl{e3d8a51b30e2072aa4bb116eb46389e342fc00fc}. The inspected sources were the research
monograph chapters

\begin{itemize}
  \item \modulepath{chapters/finite-dimensional-foundations.tex};
  \item \modulepath{chapters/projection-and-retained-spaces.tex};
  \item \modulepath{chapters/composite-band-gauge-alignment.tex}; and
  \item \modulepath{chapters/operator-comparison.tex}.
\end{itemize}

The ICMSEP compatibility bridge and the three-channel residual sign bridge use the
separate immutable owner-source snapshots recorded in Appendix~\ref{app:provenance}.
Every bridge is editorial. The inspected manuscripts supply notation and motivation;
the Formal declarations identified in Appendix~\ref{app:formal-crosswalk} determine
exact statement and proof status.

\section{ICMSEP compatibility notation}

The ICMSEP extended-abstract source pinned in Appendix~\ref{app:provenance} supplies the
overall model-adequacy problem.  Table~\ref{tab:icmsep-compatibility-bridge} records its
relation to the notation of Chapters~\ref{ch:proof-architecture} and
\ref{ch:reduction-decision}.

\begin{table}[ht]
\centering
\small
\begin{tabularx}{\textwidth}{@{}l Y Y@{}}
\toprule
Proof companion & ICMSEP source & Mathematical role \\
\midrule
$\mathfrak M_j$
  & $\mathfrak M_j$
  & candidate tight-binding model class in the declared hierarchy \\
$H_j(\theta)$
  & $\mathbf H(\boldsymbol\theta)$
  & one candidate Hamiltonian in that class \\
$\mathcal L_{E,j}(\theta)$
  & $\mathcal L_E(\boldsymbol\theta)$
  & spectral loss \\
$\mathcal L_{H,j}(C,\theta)$
  & $\mathcal L_H(\mathbf C,\boldsymbol\theta)$
  & aligned operator loss \\
$\mathcal U_j$
  & $\mathcal U$
  & admissible alignment family \\
$\mathfrak A_{E,j}$, $\mathfrak A_{H,j}$
  & $\mathfrak A_{E,j}$, $\mathfrak A_{H,j}$
  & spectral and operator admissible sets \\
$\delta_j^*$
  & $\delta_j^*$
  & metric separation of the two admissible sets \\
\bottomrule
\end{tabularx}
\caption{Compatibility-problem notation bridge to the ICMSEP extended abstract.}
\label{tab:icmsep-compatibility-bridge}
\end{table}

The source writes a minimum over $\mathbf C\in\mathcal U$.  The companion uses an
infimum until attainment has been supplied, and then recovers the minimum formulation.
It also distinguishes literal nonempty intersection from acceptance under a positive
map tolerance.  These are scope clarifications, not changes to the source's motivating
question.

\section{Finite-dimensional retained-space notation}

The source manuscript and Chapter~\ref{ch:finite-operator-foundations} use the same
representation-neutral objects with the notation shown in
Table~\ref{tab:finite-manuscript-notation-bridge}.

\begin{table}[ht]
\centering
\small
\begin{tabularx}{\textwidth}{@{}l Y Y@{}}
\toprule
Proof companion & Source manuscript & Common mathematical object \\
\midrule
$V$
  & $V$
  & $D\times M$ retained-frame embedding \\
$|v_a\rangle$
  & $|\phi_a\rangle$
  & column $a$ of the retained frame \\
$P_V=VV^\dagger$
  & $P$ or $\widehat P_{\mathcal S}$
  & ambient projector onto the retained subspace \\
$C_V(H)=V^\dagger HV$
  & $H_{\mathcal S}$
  & operator matrix in retained coordinates \\
$R_V(H)=VC_V(H)V^\dagger$
  & $PHP$ embedded in the parent space
  & reconstructed retained operator in ambient coordinates \\
$V_G=VG$
  & $V'=VG$
  & right-unitary change of retained frame \\
\bottomrule
\end{tabularx}
\caption{Representation-neutral notation bridge to the inspected source manuscript.}
\label{tab:finite-manuscript-notation-bridge}
\end{table}

The source propositions labeled \texttt{PRF-05.01} and \texttt{PRF-05.02} correspond,
respectively, to projector invariance under a retained unitary and covariance of
compression. In the Formal snapshot these are
\decl{BG-ALIGNMENT-THM-004} and \decl{BG-ALIGNMENT-THM-005}. The nearby projector and
reconstruction claims are supported by \decl{BG-ALIGNMENT-THM-001--010} under the exact
premises listed in Appendix~\ref{app:formal-crosswalk}.

\section{Bloch-frame and Fourier notation}

The source composite-band chapter writes its momentum-dependent frame as
\[
  F(\kappa)=\begin{bmatrix}u_1(\kappa)&\cdots&u_M(\kappa)\end{bmatrix}.
\]
At corresponding momentum coordinates, the notation of the proof companion is
\[
  V(k)=F(\kappa)
  =\sum_{a=1}^M |u_a(\kappa)\rangle\,{}_{\mathrm{ret}}\!\langle a|.
\]
Thus $V(k)$ and $F(\kappa)$ are two names for the map that inserts retained-coordinate
vectors into an ambient Bloch fiber.

\begin{table}[ht]
\centering
\small
\begin{tabularx}{\textwidth}{@{}l Y Y@{}}
\toprule
Proof companion & Source manuscript & Common mathematical object \\
\midrule
$V(k)$
  & $F(\kappa)$
  & momentum-dependent retained-frame embedding \\
$|v_a(k)\rangle$
  & $u_a(\kappa)$
  & column $a$ of the composite Bloch frame \\
$P_V(k)$
  & $P(\kappa)=F(\kappa)F(\kappa)^\dagger$
  & ambient projector field \\
$C_{V(k)}(H(k))$
  & $H_F(\kappa)$
  & retained operator field \\
$T_R[h]$
  & $H_U(R)$ or $\langle w_{m0}|\widehat H|w_{nR}\rangle$
  & real-space Fourier block of the retained operator \\
\bottomrule
\end{tabularx}
\caption{Bloch and Fourier notation bridge to the inspected source manuscript.}
\label{tab:bloch-manuscript-notation-bridge}
\end{table}

The companion uses $G(k)$ for an internal unitary change of one retained frame and
reserves $U(k)$ for an identification between reference and candidate retained spaces.
The source composite-band chapter uses $U(\kappa)$ for its internal gauge. These symbols
must not be identified solely because their matrices are square and unitary: an internal
frame gauge and a two-model identification have different domains, provenance, and
adequacy obligations.

\section{Aligned-difference sign bridge}

When the source manuscript's pristine operator is identified with
$h_{\mathrm{ref}}$, its doped operator with $h_{\mathrm{cand}}$, and its map $U$ runs
from reference to candidate coordinates, the source writes
\[
  \Delta H_{d\to b}=U^\dagger h_{\mathrm{cand}}U-h_{\mathrm{ref}}.
\]
The convention adopted in this companion is
\[
  \Delta_U=h_{\mathrm{ref}}-U^{-1}h_{\mathrm{cand}}U.
\]
For a unitary $U$ with the corresponding direction, this gives
\[
  \Delta_U=-\Delta H_{d\to b}.
\]
Norms are insensitive to the sign, but signed perturbations, impurity potentials, and
response coefficients are not. Any editorial extraction must preserve both the operand
order and the direction assigned to $U$.

\section{Three-channel residual sign bridge}

Section~\ref{sec:three-channel-residual-calculus} and the Lean declarations use
reference-minus-approximation conventions:
\[
  \Delta_{\mathrm{int}}^{\mathrm F}=H_{\mathrm{ref}}-H_{\mathrm{int}},\qquad
  \Delta_{\mathrm{trunc}}^{\mathrm F}=H_{\mathrm{int}}-H_{\mathrm{tr}},\qquad
  \Delta_{\mathrm{total}}^{\mathrm F}=H_{\mathrm{ref}}-H_{\mathrm{tr}}.
\]
The signed spectral residuals in the inspected \decl{Ragasa2026a} source use the opposite
order at every step:
\[
  \Delta_{\mathrm{int}}^{\mathrm A}=H_{\mathrm{int}}-H_{\mathrm{ref}},\qquad
  \Delta_{\mathrm{trunc}}^{\mathrm A}=H_{\mathrm{tr}}-H_{\mathrm{int}},\qquad
  \Delta_{\mathrm{total}}^{\mathrm A}=H_{\mathrm{tr}}-H_{\mathrm{ref}}.
\]
Consequently
\[
  \Delta_x^{\mathrm A}=-\Delta_x^{\mathrm F}
  \quad\text{for }x\in\{\mathrm{int},\mathrm{trunc},\mathrm{total}\}.
\]
Simultaneous sign reversal preserves the additive chain and every norm reported in this
companion.  It does not permit signed fields from the two conventions to be combined
without conversion.  The inspected \decl{Ragasa2026b} source presents its three
channels through a symmetric error functional $\mathcal E(X,Y)$; that scalar notation
suppresses operand orientation but does not create scalar additivity.

\section{Status of the bridge}

The source-manuscript entries \texttt{PRF-*} are prospective editorial propositions.
They do not own the exact statements or proof status. The bridge records where the
source notation and the formal declarations agree, while the Lean type and proof term
remain authoritative whenever wording or scope differs.
